% Incorporation-ready section. Requires amsmath and booktabs only.
% All state labels in this section are inherited labels, not canonical labels.
\section{Seven source rotations that expose the new fold}
\label{sec:seven-source-rotations235}

The reduction from 236 to 235 states comes from changing the arrangement of
the existing decoder.  The source still has 224 distinct DAG nodes and
6,455 instruction occurrences.  Seven closed reassociations create the
paths needed to merge states 74 and 193, while retaining the calibration
of the compact binary interpreter.  The final arrangement has maximum
depth 48, compared with 49 in the pinned 236-state input.

The new contribution is this concrete composition and its complete
validation.  Closed associativity and the use of a previously written
tape bit to distinguish merged states are already present in the
236-state predecessor.  The present construction applies these ideas
at new locations and repairs a global encoding constraint that defeated
several shorter, locally plausible candidates.

\subsection{The exact source transformation}

Write
\[
 \mathsf L:(A,(B,C))\longmapsto((A,B),C),\qquad
 \mathsf R:((A,B),C)\longmapsto(A,(B,C)).
\]
If the rotated source occurrence has prefix $p$, the left rotation changes
paths by
\[
 p0u\mapsto p00u,\qquad
 p10u\mapsto p01u,\qquad
 p11u\mapsto p1u;
\]
the right rotation uses the inverse substitutions.  Operations and
operands are unchanged.  An outer state $q$ and inner state $q'$ identify
the two corresponding pure decoder nodes immediately before that stage.
Each row of these nodes writes back its read symbol and moves right.
The inner node has exactly one incoming transition, from the stated
outer node.  These properties are checked on the literal table, rather
than inferred from the names of the source chunks.

The following sequence is applied in order.  Chunk sizes count expanded
instruction occurrences.  The last two columns describe the complete
source after the indicated stage.

\begin{center}
\begin{tabular}{rrrc rrr}
\toprule
Stage & $q$ & $q'$ & Rotation & $|A|,|B|,|C|$ & Depth & Zero run\\
\midrule
0 & --- & --- & --- & --- & 49 & 6\\
1 & 138 & 211 & $\mathsf R$ & $539,56,242$ & 48 & 6\\
2 & 116 & 134 & $\mathsf L$ & $56,121,125$ & 48 & 7\\
3 & 162 & 116 & $\mathsf L$ & $108,177,125$ & 48 & 6\\
4 & 199 & 227 & $\mathsf L$ & $2965,177,940$ & 49 & 7\\
5 & 227 & 197 & $\mathsf R$ & $1919,1046,177$ & 48 & 7\\
6 & 197 & 115 & $\mathsf R$ & $480,566,177$ & 48 & 7\\
7 & 115 & 229 & $\mathsf R$ & $406,160,177$ & 48 & 6\\
\bottomrule
\end{tabular}
\end{center}

Each stage has exactly one expanded occurrence of its rotated node.
The full path prefixes, every individual row edit, all chunk boundaries
and the intermediate source statistics are recorded in
\texttt{source\_rotation\_certificate.json}.  The portable builder
reconstructs these data directly from the pinned original leaf list.
Thus the source certificate does not depend on retaining a stochastic
search state or an undocumented compiler invocation.

Across the complete seven-stage composition, exactly 2,524 leaf paths
change.  The resulting 236-state basis differs from the input in 17
literal rows, involving 10 states.  Subsequent edits cancel some earlier
row changes, so these 17 rows should not be confused with the four rows
examined at each individual rotation.

\subsection{Why the source proof covers arbitrary register values}

For a source chunk $T$, the builder calculates a boundary signature
$\partial T=(b_T,c_T,h_T)$: upper bounds on escaping break and continue
depths, together with a flag for a halt occurrence.  For leaves it uses
\[
\begin{aligned}
 \partial\operatorname{CI}(r)&=(0,0,\mathrm{false}),&
 \partial\operatorname{CD}(r)&=(1,0,\mathrm{false}),\\
 \partial\operatorname{CB}(k)&=(k,0,\mathrm{false}),&
 \partial\operatorname{CC}(k)&=(0,k,\mathrm{false}),\\
 \partial\operatorname{CH}&=(0,0,\mathrm{true}).
\end{aligned}
\]
Writing $x^+=\max(x,0)$, a binary node satisfies
\[
 \partial(A,B)=
 \bigl(\max((b_A-1)^+,b_B),\,
       \max(c_A,(c_B-1)^+),\,h_A\lor h_B\bigr).
\]
Every one of the 21 rotated chunks has boundary
$(0,0,\mathrm{false})$.  No control action escapes a chunk, so either
parenthesization executes the same closed chunks in the same order.
This argument includes divergence inside a chunk; it does not require
the chunks to terminate.

The executable certificate also checks the complete guarded source
graph at every stage.  It numbers leaves in their common left-to-right
order and records the operation, register operand, and every successor.
For a decrement it checks both the zero branch and the nonzero branch.
Break and continue destinations are resolved against the entire source
path, including enclosing contexts.  The initial leaf is checked as
well.  All 6,455 ordered operations and all 8,386 guarded cases agree
with the original source.  Successor resolution also checks that any
unused suffix contains only zeros, as required by the source interface.

At corresponding leaves, the two programs therefore apply exactly the
same register operation and select the same next leaf for every
natural-valued register file.  Induction gives identical finite source
runs; the same step correspondence preserves infinite runs.  This is
an unbounded source equivalence proved by a finite graph comparison,
not a test on a finite range of counter values.  A separate checker
using a trie and ancestor unwinding verifies the same source graph
without importing the builder's arithmetic successor resolver.

\subsection{What the four core rotations create}

Before the final fold, the relevant native rows are
\[
\begin{array}{c|cc}
 &0&1\\\hline
 74&0\ R\ 114&1\ R\ 233\\
 193&0\ R\ 190&1\ R\ 213.
\end{array}
\]
The first three rotations construct two useful binary routing patterns:
\[
\begin{array}{c|cc}
 &0&1\\\hline
 138&0\ R\ 210&1\ R\ 211\\
 116&0\ R\ 246&1\ R\ 134.
\end{array}
\]
Existing or newly arranged following rows supply
\[
 210/0\to114,\quad211/0\to190,\quad
 246/1\to233,\quad134/1\to213,
\]
each preserving its read symbol and moving right.  The fourth rotation
makes the old 193 entry arrive immediately after a written 1.  The final
three repairs preserve that property while moving its source occurrence.
In the final basis, every arrival at 74 is through $140/0=0\ R\ 74$,
while the only arrivals at 193 are $163/1=1\ R\ 193$ and
$229/1=1\ R\ 193$.  State 74 consequently has a 0 immediately to the
left of the head; state 193 has a 1 there.

Redirect both incoming 193 targets to 74, delete state 193, and set
\[
 74/0=0\ L\ 138,\qquad 74/1=1\ L\ 116.
\]
The state records the current bit by choosing 138 or 116, moves left to
read the old-state tag, and follows a pre-existing rightward path back
through the original cell.  Each of the four exceptional old rows is
simulated in exactly three transitions.  The old and new endpoints have
the same state, head position one cell to the right, and entire tape.
All other 468 old rows are simulated in one transition after renaming
193 to 74.  This is the one-state saving; the rotations themselves
neither add nor remove a working state.

Relative to the original 236-state table, the complete result changes
20 surviving rows and removes the two rows of state 193.  The retained
state labels are inherited labels; the separately supplied canonical
table uses working states $0,\ldots,234$ and halt outcome 235.

\subsection{Why the last three rotations are included}

The four core rotations already give the source-graph proof and the
local state-fold proof.  They do not yield a correct binary interpreter.
They create a source path containing seven consecutive zeros, conflicting
with the compact tape encoding.  In the actual failed candidate, the
universal runtime check rejects the increment at leaf ordinal 648,
whose path is
\begin{center}\ttfamily
1110100100110110101010000000.
\end{center}
Here and below leaf ordinals start at zero.  A concrete finite execution
independently exposes the same failure.

The final three rotations transport a 177-leaf closed chunk through
successive larger chunks, restoring a maximum zero run of six.  The
repair search allowed intermediate source representations with seven
or eight consecutive zeros; it demanded the six-zero condition at the
final binary candidate.  All intermediates remain source-equivalent.
They are algebraic stages of the construction, not individually
calibrated binary machines.  Enforcing every final encoding constraint
after every source move would exclude the displayed route.

Restoring the zero bound alone is insufficient.  One rejected repair
rotated the decoder pair 105/123.  Although the guarded source graph
still agreed, that edit changed bits immediately preceding a shared
arithmetic routine.  A backward-reading helper then took the wrong
branch for an increment at leaf 195.  It also changed the exact blank
startup at step 12,742.  The final construction avoids that edit.
This failure is useful evidence about the optimization problem:
decoder nodes can simultaneously serve as runtime helper states, and
their behavior can depend on preceding path bits.

These observations motivate three distinct checks.  Closed-source
reasoning proves the intended register program is unchanged.  The
all-tape local fold proof proves the final deletion implements the
rotated literal basis.  The instruction-runtime proof establishes
that this exact basis actually executes the transformed source.  None
of these obligations can be replaced by the other two.

\subsection{The final binary checks and reproducible construction}

The seven-rotation 236-state basis and the final 235-state machine both
pass fresh unbounded runtime validation on all 8,386 guarded cases.
The runs use 2,266,171 and 2,269,055 symbolic rules respectively, cover
all 472 and 470 literal rows, and establish minimum instruction duration
13.  The runtime contracts cover arbitrary natural values of all nine
primary registers, arbitrary nonnegative extra gap, arbitrary first
spare register, and an arbitrary finite spare tail.  They check the
complete tape, head and next source location.

The blank startup is unchanged: all 409,836 checkpoints, including
the initial one, agree through 409,835 transitions.  No changed or
deleted row is visited during this startup.  The endpoint has inherited
dispatcher state 96, head position $-5$, and exactly 351 one-cells,
with their full support recorded in the startup certificate.

The standard-library Python builder accepts the extracted pinned
236-state package through \texttt{--baseline-dir}, and an output
directory through \texttt{--output-dir}.  It verifies both exact input
hashes, replays the seven rotations, checks every guarded source graph,
builds and renumbers the final table, and replays the complete startup.
The independent source, fold and universal runtime checkers are
separate proof artifacts.  The construction is deterministic and
requires no discovery search, network access or original compact
compiler.

The construction and these finite proof certificates compose with
the pinned logical calibration as ordinary mathematics formalizable
in ZFC.  They do not by themselves claim a newly compiled proof-assistant
endpoint for 235 states.  The discovery search statistics and finite
execution filters are useful provenance, but are not premises of the
calibration theorem.
